As the first self-referential agent, \framework has to construct all task-related code autonomously, which poses significant challenges. Consequently, this work does not compare directly with the most complex existing agent systems, such as OpenDevin~\citep{wang2024opendevinopenplatformai}, which have benefited from extensive manual engineering efforts. 
% Currently, \framework is not sufficiently stable and may be prone to error accumulation, hindering its ability to continue self-optimization. 
This makes it unrealistic to expect it to outperform systems that have taken researchers several months or even years to develop. The experiments presented in this paper are intended to demonstrate the feasibility of recursive self-improvement. 

% The experiments presented in this paper are intended to demonstrate the effectiveness and feasibility of recursive self-improvement. 

Additionally, as the agent system becomes increasingly complex through self-optimization, it may require exponentially more intelligence to understand itself. 
Consequently, a system capable of complete self-referential at the outset may lose this capability as it evolves~\citep{Yampolskiylimitation}. The exact point at which the agent can no longer comprehend and improve itself has not been thoroughly explored. Investigating this phenomenon, both experimentally and theoretically, could provide valuable insights into the limitations of recursive self-improvement.
A more robust and advanced implementation of the \framework is anticipated, with numerous potential improvements outlined in Section \ref{sec:future}. 